% Generated from validated Article Draft Result. Do not hand-edit; revise the structured draft instead.
\section{知識を重みに閉じ込めない――REALMからRAGへ}
\label{pkg:retrieval-grounding}
\noindent\textbf{大規模modelの知識をすべてparameterへ埋め込む以外に、必要なdocumentを外部corpusから検索してmodelへ渡すという設計が2020年に明確になった。REALMはretrievalをpretrainingへ組み込み、RAGはretrieverとgeneratorをknowledge-intensive taskで接続した。retrievalは『検索機能の追加』ではなく、parametric memoryの外側に知識surfaceを置くarchitectureへ向かう分岐だった。}\autocite{src-38c08ba169f1,src-4bf65d8fc7dc,src-c441b6f2b8f4}

REALMが問題にしたのは、language modelのworld knowledgeをparameterだけに保存する設計である。論文はlatent knowledge retrieverをpretrainingへ組み込み、大規模document corpusから関連documentを検索してattendする仕組みを示した。retrievalはfine-tuning後の補助処理ではなく、pretraining・fine-tuning・inferenceを通じて使われるknowledge access mechanismとして位置付けられていた。\autocite{src-4bf65d8fc7dc}


RAGはretrievalとgenerationの接続をさらに明示した。retrieverが外部corpusからdocumentを取得し、generatorがそのretrieved contextを条件としてoutputを生成する構造を取り、knowledge-intensive NLP taskを対象に評価された。ここではretrievalが単にanswer候補を返すのではなく、生成modelのconditioning sourceとして組み込まれている点が重要である。\autocite{src-38c08ba169f1,src-c441b6f2b8f4}


REALMからRAGへの流れを年次trajectoryとして見ると、知識をmodel parameterへ固定するparametric memoryと、外部corpusを必要時に参照するnon-parametric knowledge surfaceが分離し始めたことが分かる。これは後年のgrounded generationをそのまま2020年へ遡及させる話ではない。2020年に確認できるのは、retrievalをmodel architectureの一部として学習・生成へ接続する設計が主要研究課題になったことである。\autocite{src-38c08ba169f1,src-4bf65d8fc7dc,src-c441b6f2b8f4}


\begin{claimboundary}
REALMとRAGのbenchmark結果、retrieval accuracy、generation qualityなどの評価は各author / projectの報告である。本稿はretrieval-augmented pretrainingとretrieval-conditioned generationという設計上の位置付けを一次資料から確認しているが、性能優位や知識更新能力を独立再現したものではない。\autocite{src-38c08ba169f1,src-4bf65d8fc7dc,src-c441b6f2b8f4}
\end{claimboundary}
